%% kapa: exam=1 round=31 year=2020 period=1 form=A subject=economics subject_ko=경제학원론 seq=12 number=52
%% source: Q-Net 감정평가사 기출문제 (cst00309) artlSeq=5207219 file=extracted/1008/5207219/A형/A형/2020년도 제31회 감정평가사 1차 시험_1교시_A형.pdf
%% conversion: kordoc 4.24.51 -> tools/gen1.py (auto); review=proofread-sonnet; flags=none
\begin{question}{52}
보수 행렬이 아래와 같은 전략형 게임(strategic form game)에서 보수 a값의 변화에 따른 설명으로 옳은 것은? (단, 보수 행렬의 괄호 안 첫 번째 값은 甲의 보수, 두 번째 값은 乙의 보수이다.) \par
{\small\centering
\begin{tabular}{|c|c|c|c|}
\hline
\multicolumn{2}{|c|}{} & \multicolumn{2}{c|}{乙} \\ \cline{3-4}
\multicolumn{2}{|c|}{} & 인상 & 인하 \\ \hline
\multirow{2}{*}{甲} & 인상 & (a, a) & (-5, 5) \\ \cline{2-4}
 & 인하 & (5, -5) & (-1, -1) \\ \hline
\end{tabular}\par}
\begin{choices}
  \item a > 5 이면, (인상, 인상)이 유일한 내쉬균형이다.
  \item -1 < a < 5 이면, 인상은 甲의 우월전략이다.
  \item a < -5 이면, 내쉬균형이 두 개 존재한다.
  \item a < 5 이면, (인하, 인하)가 유일한 내쉬균형이다.
  \item a = 5인 경우와 a < 5인 경우의 내쉬균형은 동일하다.
\end{choices}
\end{question}
\answer{4}
